\section{Broader Impacts}
\label{sec:broader}

\noindent \textit{\textbf{Societal and scientific impact.}} Improving the reliability of computer systems has critical importance for the \textit{nation's technology-driven industry and economic competitiveness}. If successful, our proposed techniques and associated systems can dramatically improve software quality by eliminating bugs that cause massive financial losses~\cite{blackout-bug,knights-bug-cost,heartbleed-cost}, theft of private information~\cite{equifax-breach}, and security breaches of large scale Internet infrastructure~\cite{mirai}. This project will also reduce the exorbitant costs associated with finding, reproducing and fixing bugs (\textasciitilde300 billion dollars~\cite{goodliffe_2015_BecomingBetterProgrammer,tombritton_2013_ReversibleDebuggingSoftware,seacord_2003_ModernizingLegacySystems}). The outcomes of the successful completion of this project will be (1) a real-world bug reproduction system with unprecedented effectiveness, accuracy, and efficiency, (2) a large number of new concepts that advance the state of the art in program analysis and system support for bug reproduction (trace-driven and adaptive symbolic execution, hybrid slicing, schedule refinement via adaptive monitoring, hybrid concurrency bug detection, and schedule synthesis). Finally, the developed techniques will benefit other areas such as software testing, debugging, and program repair.


\noindent \textit{\textbf{Dissemination of results.}}
We will publish our results at premier computer systems conferences and accompany the papers with well-documented software releases with permissive licenses (details are in the data management plan) as we have done in the past~\cite{racemob,gist,huron}. Reproducibility is a central theme of this proposal, both in terms of the research pursued, the dissemination of results, and as my career goal. I will release all the tools, datasets and scripts allowing to recreate the results for promoting adoption and further studies. \textbf{\textit{I am also co-leading the first artifact evaluation track~\cite{sysartifacts} in a major systems conference~(SOSP'19)}}, where we will adopt an evaluation model from the programming languages community that has well-established processes for reproducibility~\cite{pldiae}. \textit{I am fully committed to making artifact evaluation mainstream in my community}. Finally, my recent findings on microprocessor vulnerabilities (\ie Foreshadow~\cite{foreshadow,foreshadow-ng}) have been widely featured in popular press (see {\small \url{http://foreshadowattack.com}}); I will explore similar widespread dissemination for this project. 

\noindent \textit{\textbf{Broadening participation.}} I am committed to the \textit{development of a diverse, globally competitive STEM workforce}. The results of this research will inform my mentorship of “GirlsEncoded”~\cite{girlsencoded} and “Explore Graduate Studies”~\cite{ExploreG5:online}, which are undergraduate outreach programs at the University of Michigan, as well as my involvement of undergraduates in computer systems research. To date, 17 undergraduates (\textit{8 women}) have fully participated in my research, which led to top-tier papers (in OSDI'18~\cite{rept} and PLDI'19~\cite{huron}). Some of these students have joined top graduate programs such as Stanford, CMU, UIUC, UT Austin, etc. Funding for these students will be obtained from special programs that are available at NSF or through Michigan's College of Engineering programs: summer undergraduate research experience~\cite{sure} and undergraduate research opportunity program~\cite{urop}. Finally, the educational plan of this proposal will develop a course centered around teaching problem-solving skills through debugging that will be piloted in the Qualcomm Thinkabit lab in the Michigan Engineering Zone in Detroit, which provides a \textit{free and accessible STEM experience for underrepresented youth} in the Southeast Michigan area.

\noindent \textit{\textbf{Increased partnerships between academia and industry.}} We will leverage our collaborations with industry researchers and practitioners, including Microsoft and Intel (see attached letters), to foster technology transfer. For instance, my Gist~\cite{gist} and RaceMob~\cite{racemob} projects have been used by different research groups in academia and industry (e.g., at Intel and Microsoft). My collaboration with Microsoft resulted in a dynamic instrumentation framework~\cite{bfs} used at Microsoft. Similarly, REPT~\cite{rept}, which I co-developed, is used in hundreds of millions of real-world systems and is an integral part of the Windows debugger~\cite{windbg}. 
% established players like ARM where the PIs have existing collaborations. 
